%\chapter*{附錄：常見的問題}
%\addcontentsline{toc}{chapter}{附錄：常見的問題}

\appendix
\renewcommand{\appendixname}{附錄}
\renewcommand{\thechapter}{\arabic{chapter}}
\renewcommand{\thefigure}{\appendixname~\thechapter.\arabic{figure}}  % Appendix 1.1, 1.2, …

\refstepcounter{chapter}
\chapter*{附錄：常見的問題}
\addtocontents{toc}{\protect\setlength{\cftbeforechapskip}{0.3em}}
\addcontentsline{toc}{chapter}%
  {\appendixname~\thechapter\quad 常見的問題}



\section*{當你深愛的弟兄（或是姐妹）離世了，可是你走不出傷痛。}

你必定很深地想念那位與你一同走過生命的弟兄或是姐妹，也很深地渴望天家。我想，主知道你心裡的思念，也知道這段離別的路有多麼不容易。但今天，主仍然把你留在世上。你不是被留下來獨自承受痛苦，而是被主的恩典追上的人，將來還要被主的恩典帶入榮耀裡。

有一天，我們要與天地的主宰基督同在。而我們最大的盼望，就是永遠與主同在，親眼見祂的榮耀。在祂裡面，我們也有與所有屬主的人一同進入永恆的盼望。

願姊妹在思念弟兄的日子裡，仍然仰望基督。祂不僅是我們今生的倚靠，也是我們永恆的盼望。願主親自扶持你，使你在每一天的生活中，都能經歷祂的同在與恩典。
主給我們的應許，是祂要擦去我們的眼淚。

傳道書 3:1、4 告訴我們，凡事都有定期，哭有時，笑有時；哀慟有時，跳舞有時。
如果你開始快樂，不是代表沒有那麼愛他了。如果你去旅行、或參加聚會，也不是代表你已經忘記了他了。
你可以繼續珍惜與先生共同生活的回憶，也可以慢慢學習適應沒有先生陪伴的日子。


\section*{當復活的時候，人也不娶也不嫁，乃像天上的使者一樣（馬太福音 22:30）}

我只能想像啟示錄 7:9--10 所描繪的那個場景：
「有許多的人，沒有人能數過來，是從各國各族各民各方來的，站在寶座和羔羊面前，身穿白衣，手拿棕樹枝，大聲喊著說：『願救恩歸與坐在寶座上我們的神，也歸與羔羊。』」

也許，在那壯闊的隊伍中，你看見了自己的配偶。
可是，你眼前的他，已經不再是你在地上所熟悉的模樣。他身穿白衣，眼中有光，正興奮地讚美救恩的偉大。他全心全意地敬拜那位救贖他的主。此刻，他的生命完全以基督為中心；他所渴慕的，是在永生裡親近主、認識主、讚美主，與主同行。

你也許會發現，將來的他，遠超過你今生所認識的他。
你從前很了解這位弟兄。你知道他的事業曾經是他生活中的重心，也知道他每天上班前習慣喝一杯咖啡。你熟悉他的習慣、他的喜怒哀樂，也陪伴他走過人生的高低起伏。然而，到了那日，他將不再受今世的軟弱、勞苦與憂愁所限制，而是得著不朽壞、榮耀的身體（哥林多前書 15:42--44）。

可是，你也變了。
你同樣站在那群從各國各族各民各方而來的人中間，身穿白衣，手拿棕樹枝，向坐在寶座上的神和羔羊獻上讚美。那位最完全地認識你、最深地愛你，也為你捨命的救主，如今就在你面前。你將以全新的生命愛祂、敬拜祂，在永恆中與祂同行。

哥林多前書 13:12 說：「我們如今彷彿對著鏡子觀看，模糊不清……到那時就要面對面了。」我們今天對主的認識仍然有限，但將來要更完全地認識祂。到了那日，我們不再只是憑著信心仰望主，而是要親自見到祂。我們將與眾聖徒一同站在羔羊面前，敬拜那位為我們捨命、救贖我們的主，永遠享受祂的同在。



\section*{如果將來我在天國，所愛的家人卻不在那裡，我還能快樂嗎？}

當我們想到自己的父母、配偶、兒女或兄弟姊妹，如果有一天我們在天國，而所愛的人卻不在那裡，我們怎麼可能不感到悲傷呢？

我們也許會問：到了永恆裡，我們是否仍然記得今生所愛的人？如果記得，我們是否仍然愛他們？如果愛仍然存在，而他們卻沒有得救，我們又怎麼可能在沒有悲傷的情況下享受永恆？

聖經確實啟示了永恆的盼望、神的公義，以及祂要除去一切悲哀的應許；但聖經並沒有完整解釋，得救者在新天新地中將如何經歷對失喪親人的記憶與感情。因此，我們需要謙卑承認，聖經沒有明確說明的部分，不應被我們的推測取代。

第一點，到了天國，我們是否還會記得不信主的家人？我們知道耶穌復活以後，祂仍然是那位被釘十字架、後來復活的主。神也會改變形狀我們這卑賤的身體，使我們和他自己榮耀的身體相似。（腓立比書3:21）

第二點：如果仍然記得，看到所愛的人不在天國，怎麼可能快樂？
假設一位母親在天國，卻知道自己的兒子沒有得救；或者一位丈夫在神的同在中，卻知道自己所愛的妻子沒有進入永恆的生命。我們很難想像，他怎麼可能完全沒有悲傷。

如果愛仍然存在，失去所愛之人的痛苦怎麼可能消失？
啟示錄21:4應許我們：
「神要擦去他們一切的眼淚；不再有死亡，也不再有悲哀、哭號、疼痛，因為以前的事都過去了。」
神要親自擦去祂子民一切的眼淚，使死亡、悲哀、哭號與疼痛不再成為他們生命的一部分。

然而，這段經文並沒有詳細說明，神將如何更新我們的記憶、情感與理解，使我們在新天新地中完全享受祂的同在。
我們所能確信的是，神的應許不會落空。

此外，到了永恆，我們將完全承認神的審判是公義的。啟示錄 16:7 宣告：
「主神啊，你的判斷義哉！誠哉！」
祂的判斷完全公義，也完全符合祂自己的聖潔。
然而，承認神的審判完全公義，並不等於我們已經知道，得救者將如何理解或經歷未得救親人的結局。

第三點：神的聖潔與憐憫並不矛盾。
耶穌在地上時，祂知道人的罪，也知道審判的真實，卻仍然為耶路撒冷哀哭。
馬太福音 23:37 記載，耶穌說：
「耶路撒冷啊，耶路撒冷啊……我多次願意聚集你的兒女，好像母雞把小雞聚集在翅膀底下，只是你們不願意。」
這段經文讓我們看見，耶穌的聖潔與祂對人的憐憫並不矛盾。
祂沒有因為知道人的罪，就失去對人的憐憫；祂也沒有因為愛人，就否定神的真理與審判。祂的憐憫是真實的，祂的審判也是公義的。

那麼，今天的我們應當如何面對這個問題？
我們仍然可以關心家人的屬靈生命，為他們禱告，也可以為他們流淚。
保羅曾為以色列同胞深深憂傷，甚至說自己為他們大有憂愁，心裡時常傷痛（羅馬書9:2-3）。
這提醒我們，對所愛之人的得救有深切的負擔，並不表示我們缺乏信心。
今天，我們仍然可以為所愛的人禱告，向他們分享福音，也將自己無法掌握的事交託給神。